\sectiontitle{一、单项选择题}{10}{20}

% C01 S02:16 | 基础
\mc{同一初始速度误差下，两组比例增益试验结果如下：$K_p=0.4$ 时，稳定时间为 $1.5$ s、残余误差为 $0.6$ m/s，未见振荡；$K_p=0.8$ 时，稳定时间为 $0.9$ s、残余误差为 $0.2$ m/s，但有持续振荡。哪项结论符合这些数据？}
{较大的 $K_p$ 缩短了稳定时间并减小残余误差，但振荡提示稳定性风险增加}
{较大的 $K_p$ 已完全消除稳态误差，所以不需要再检查振荡}
{较小的 $K_p$ 响应更快，因为其稳定时间更长}
{两组结果相同，因为比例项不影响闭环响应}
{A}{数据表明增大 $K_p$ 可加快响应并减小稳态误差，但也可能增加振荡风险。}

% C02 S09:14 | 基础
\mc{课件中的 Q 表在状态 $S_3$ 对动作 $A_1,A_2,A_3,A_4$ 的值依次为 $-1,1,0,-2$。若当前只按最大 Q 值选动作，应选择：}
{$A_1$，因为它的 Q 值绝对值最大}
{$A_2$，因为它的 Q 值最大}
{$A_3$，因为它的 Q 值最接近零}
{$A_4$，因为负值表示奖励更高}
{B}{比较该行数值，最大值为 $1$，对应 $A_2$。最大绝对值不等于最大 Q 值。}

% C03 S10:10 | 基础
\mc{车辆以 $72$ km/h 行驶，GNSS 每秒输出 $10$ 次位置。按相邻输出间匀速估算，车辆每次更新间隔移动多少米？}
{$0.2$ m}
{$1$ m}
{$2$ m}
{$20$ m}
{C}{先换算速度：$72$ km/h $=20$ m/s。更新间隔为 $0.1$ s，因此距离为 $20\times0.1=2$ m。}

% C04 S03:6 | 基础
\mc{道路参考边界的定位误差需小于 $0.2$ m。可选地图的标称精度分别为 $0.03$ m 和 $1.0$ m。哪项选择与该误差要求相符？}
{两种地图都相符，因为二者都含有道路信息}
{只用 $1.0$ m 地图，因为它覆盖范围更大}
{先把 $1.0$ m 精度换成厘米，就会优于 $0.03$ m}
{选 $0.03$ m 精度地图；$1.0$ m 精度不足以满足给定误差要求}
{D}{给定误差阈值为 $0.2$ m。$0.03<0.2$，而 $1.0>0.2$。}

% C05 S04:16-17,25-26 | 基础
\mc{一张车道掩码已标出所有车道线像素，但还要区分左车道线与右车道线。应增加哪一步？}
{把掩码再按固定阈值转成一次二值图}
{使用像素 embedding 并按嵌入聚类区分实例}
{只判断整张图是否存在车道线}
{只对每个像素重新计算灰度值}
{B}{二值语义掩码只标出车道像素。像素 embedding 与聚类用于区分不同车道实例。}

% C06 S06:9,11 | 基础
\mc{路口地图允许车辆所在车道通向左转或直行车道，历史运动也支持两种意图。哪种预测结果更适合后续规划？}
{保留左转与直行候选轨迹，并记录各自的可能性}
{把两种轨迹逐点平均，并将平均轨迹当作唯一真实未来}
{只复制当前检测框，令目标未来位置不变}
{只保留地图中几何长度较短的车道}
{A}{地图连接关系和历史运动可支持多个未来意图。规划需要保留这些可能性，不能把平均轨迹视为真实轨迹。}

% C07 S02:18 | 基础
\mc{车速测量含有快速抖动。若微分增益较大，哪项现象最需要检查？}
{误差历史被积分后造成长期累积}
{快速测量变化经微分作用放大，导致控制量抖动}
{稳态误差必定被微分项完全消除}
{控制器自动忽略所有高频测量噪声}
{B}{微分项响应误差变化率，对高频噪声敏感。应检查控制输出抖动，并谨慎设定微分作用。}

% C08 S05:29 | 基础
\mc{半径离群点滤波设置邻域半径 $r$，保留条件为邻居数不少于 $k^*=3$。某点的半径邻域内只有 $1$ 个邻居，应如何处理？}
{保留，因为该点自身也计入邻居数}
{保留，因为半径滤波只按点到原点的距离判断}
{删除，因为邻居数 $1<3$}
{删除，因为该点的反射强度未给出}
{C}{题目给出邻居数为 $1$，低于保留阈值 $3$，因此该点属于孤立点并删除。}

% C09 S03:17-19 | 基础
\mc{OpenDRIVE 参考线上的两个位置分别为 $s=12$ m 和 $s=17$ m。后一个位置相对前一个位置沿参考线前进多少米？}
{$-5$ m}
{$0$ m}
{$5$ m}
{$29$ m}
{C}{沿参考线的位置差为 $17-12=5$ m。}

% C10 S09:16 | 基础
\mc{训练样本按时间顺序连续送入网络时，连续样本高度相关。哪种经验回放做法可打散训练样本的时间顺序并重用旧样本？}
{只保留最新一步的奖励}
{每轮只重复同一个动作，不保存状态}
{删除下一状态，只保留网络参数}
{从经验池中随机抽取转移样本组成批次}
{D}{经验回放保存状态转移，并随机抽样训练。它能提高样本利用率并减弱连续样本的顺序相关性。}

\exampagebreak
\sectiontitle{二、填空题}{10}{20}
\textbf{说明：}每题一个空，每空 2 分。数值结果按题目给出的单位作答。

% C11 S02:19 | 中等
\fillq{位置式 PID 定义为 $U(k)=K_pe(k)+K_I\sum_{i=0}^{k}e(i)+K_D[e(k)-e(k-1)]$。已知 $e(0)=1$、$e(1)=1$、$e(2)=2$，$K_p=0.5$、$K_I=0.2$、$K_D=0.1$，则 $U(2)=$\blankline。}
{$1.9$}{代入得 $0.5\times2+0.2\times(1+1+2)+0.1\times(2-1)=1+0.8+0.1=1.9$。}

% C12 S02:19 | 中等
\fillq{增量式 PID 中，$\Delta U(k)=K_p[e(k)-e(k-1)]+K_Ie(k)+K_D[e(k)-2e(k-1)+e(k-2)]$。已知 $e(k)=2$、$e(k-1)=1$、$e(k-2)=0$，$K_p=0.4$、$K_I=0.2$、$K_D=0.1$，则 $\Delta U(k)=$\blankline。}
{$0.8$}{比例项为 $0.4$，积分项为 $0.4$，二阶差分项为 $0$，所以增量为 $0.8$。}

% C13 S09:14 | 中等
\fillq{按 $Q_{new}=Q+\alpha[r+\gamma M-Q]$ 更新。已知 $Q=2$、$\alpha=0.25$、$r=-1$、$\gamma=0.8$、$M=5$，则 $Q_{new}=$\blankline。}
{$2.25$}{目标为 $r+\gamma M=-1+0.8\times5=3$。因此 $Q_{new}=2+0.25(3-2)=2.25$。}

% C14 S10:23 | 中等
\fillq{直方图滤波的先验概率为 $[0.2,0.5,0.3]$，观测似然为 $[0.2,0.6,0.2]$。逐项相乘并归一化后，后验概率为 \blankline。}
{$[0.1,0.75,0.15]$}{未归一化概率为 $[0.04,0.30,0.06]$，总和为 $0.40$。除以总和得到 $[0.1,0.75,0.15]$。}

% C15 S10:24 | 中等
\fillq{四个位置按环形排列，概率为 $[0.1,0.2,0.3,0.4]$。位置确定地向右移动一格，且末位置后回到首位置。移动后的概率为 \blankline。}
{$[0.4,0.1,0.2,0.3]$}{每个位置的概率移动到右侧相邻格。末位置的 $0.4$ 回到首位置。}

% C16 S09:12,14 | 中等
\fillq{折扣回报定义为 $R_t=r_t+\gamma r_{t+1}+\gamma^2r_{t+2}$。若奖励依次为 $3,2,-1$ 且 $\gamma=0.5$，则 $R_t=$\blankline。}
{$3.75$}{计算得 $3+0.5\times2+0.25\times(-1)=3.75$。}

% C17 S03:6 | 中等
\fillq{导航地图标称精度为 $1.2$ m，高精度地图标称精度为 $0.03$ m。前者的误差尺度约为后者的 \blankline 倍。}
{$40$}{比值为 $1.2/0.03=40$。这是标称精度尺度之比，不表示每个位置都产生恰好 40 倍误差。}

% C18 S07:23 | 中等
\fillq{两时刻的预测位置为 $(1,0)$ m、$(2,2)$ m，真实位置为 $(0,0)$ m、$(2,0)$ m。定义平均位移误差为各时刻欧氏距离误差的算术平均，则该误差为 \blankline。}
{$1.5$ m}{两个时刻的误差分别为 $1$ m 和 $2$ m，平均值为 $(1+2)/2=1.5$ m。}

% C19 S05:35 | 中等
\fillq{查询点为 $(0,0,0)$ m，三个候选点为 $(2,1,2)$ m、$(-1,0,0)$ m 和 $(0,0,2)$ m。按三维欧氏距离，最近邻距离为 \blankline。}
{$1$ m}{三点距离分别为 $3$ m、$1$ m 和 $2$ m，最小值为 $1$ m。}

% C20 S10:9 | 中等
\fillq{按课件的 RTK 近距离误差近似条件，基站与车载接收机相距必须小于 $30$ km。候选基站距离为 $26$ km 和 $32$ km，应选距离为 \blankline 的基站。}
{$26$ km}{只有 $26<30$ km 满足给定的距离条件。距离条件本身不能替代基站、卫星和环境质量检查。}

\exampagebreak
\sectiontitle{三、问答题}{3}{30}

% C21 S10:23-26 | 中等
\essay{某车在无 GNSS 环境中使用四格直方图表示位置概率。先验为 $[0.4,0.3,0.2,0.1]$，位置编号为 $0,1,2,3$。车辆确定地向右移动一格，且位置按环形相连。随后观测的似然为 $[0.1,0.7,0.1,0.1]$。
\begin{subquestions}
\item 计算移动后的预测概率。（3 分）
\item 将预测概率与观测似然逐项相乘并归一化，写出后验概率。（4 分）
\item 指出最大后验位置，并说明为何后续仍需新的观测或其他定位信息。（3 分）
\end{subquestions}}
{\textbf{参考答案与评分：}
\begin{subquestions}
\item 概率向右循环移动一格，得 $[0.1,0.4,0.3,0.2]$，得 3 分。
\item 未归一化值为 $[0.01,0.28,0.03,0.02]$，得 2 分；总和为 $0.34$，归一化为 $[1/34,28/34,3/34,2/34]$，得 2 分。
\item 最大概率位于位置 1，得 1 分。运动预测会传播或增加位置不确定性，观测可补充信息；需要持续更新，得 2 分。
\end{subquestions}}

% C22 S03:9;S07:19-23,27 | 综合
\essay{自车接近无信号交叉口。地图显示两条可行车道。感知到一组行人后，预测模块给出两种群组运动模式：左转概率为 $0.6$，右转概率为 $0.4$。若自车立即通过，左转模式下的碰撞代价为 $100$，右转模式下代价为 $0$；若自车等待，两个模式下的代价均为 $5$。
\begin{subquestions}
\item 计算立即通过的期望代价，并与等待代价比较。（3 分）
\item 说明为何应保留两种预测模式，而不应只用逐点平均轨迹规划。（3 分）
\item 提出两项检查，用于评估该预测和决策在开放环境中的可靠性。（4 分）
\end{subquestions}}
{\textbf{参考答案与评分：}
\begin{subquestions}
\item 立即通过的期望代价为 $0.6\times100+0.4\times0=60$，得 2 分；等待代价为 $5$，故按给定代价应等待，得 1 分。
\item 群组预测保留了组内、组间及个体轨迹关系，得 1 分；左转和右转是不同的运动模式，逐点平均可能不对应实际会出现的模式，得 1 分；应将候选轨迹及其可能性传给规划并检查风险，得 1 分。
\item 按群组和路口场景划分独立验证数据，检查两种模式的覆盖情况，任一项得 2 分；比较预测轨迹与实际轨迹的误差，或检查碰撞、等待决策及长尾场景，任一项得 2 分。答案须说明所检查的数据或指标。
\end{subquestions}}

% C23 S02:15-18,29;S09:4,8,13 | 综合
\essay{某次控制器调参记录包含两个比例增益试验。$K_p=0.4$ 时，响应时间为 $1.5$ s、稳态误差为 $0.6$ m/s，未见明显振荡；$K_p=0.8$ 时，响应时间为 $0.9$ s、稳态误差为 $0.2$ m/s，但有持续振荡。请结合 PID 环节分析结果。
\begin{subquestions}
\item 说明提高 $K_p$ 对本次响应时间和稳态误差的影响。（2 分）
\item 若任务要求消除剩余稳态误差，说明应调整哪个环节，并指出其对系统动态响应的一项影响。（4 分）
\item 测量信号含高频噪声时，说明为何不能仅通过增大微分作用来改善动态性能，并给出一项参数选择原则。（4 分）
\end{subquestions}}
{\textbf{参考答案与评分：}
\begin{subquestions}
\item 稳定时间从 $1.5$ s 降至 $0.9$ s，稳态误差从 $0.6$ m/s 降至 $0.2$ m/s，得 2 分。
\item 应考虑积分环节，得 1 分。课件指出积分作用可消除稳态误差，得 1 分；积分时间常数减小会增强积分作用，得 1 分；积分作用过强会降低稳定性并使动态响应变慢，得 1 分。
\item 微分环节放大噪声影响，得 2 分；微分作用过强会损害抗干扰能力，得 1 分；参数应在响应速度、稳态误差、振荡与噪声敏感性之间比较，得 1 分。
\end{subquestions}}

\exampagebreak
\sectiontitle{四、程序题}{2}{30}

% C24 S04:25-26 | 中等
\programq{LaneNet 等分割方法会输出车道线掩码。请编写 \code{lane_center(mask)}，拟合直线路段的左右车道位置。

\textbf{输入：}有效的二维二值 NumPy 数组，形状为 $(H,W)$，其中 $H,W\ge5$。非零像素表示车道线。$x$ 向右，$y$ 向下。左右车道线分别位于图像中心两侧，并近似为直线。

\begin{subquestions}
\item 仅使用 $y\ge\lfloor H/2\rfloor$ 的非零像素。令 $c=(W-1)/2$，按 $x<c$ 和 $x>c$ 分为左右组，忽略 $x=c$ 的像素。（5 分）
\item 每组至少包含两个不同的 $y$ 值，否则返回 \code{None}。分别以最小二乘法拟合 $x=ay+b$，计算底行 $y=H-1$ 的左右横坐标。（5 分）
\item 若结果不满足 $0\le x_L<x_R\le W-1$，返回 \code{None}。否则返回 \code{left_x}、\code{right_x}、\code{center_x} 和 \code{offset_px}。车道中心为 $(x_L+x_R)/2$，偏移为 $c-\text{车道中心}$。写出样例输出和偏移单位。（5 分）
\end{subquestions}
\codefile{snippets/C01.py}
\textbf{可用 API：}\code{np.nonzero}、\code{np.unique}、\code{np.polyfit(y,x,1)}。最后一个 API 返回 $(a,b)$。边界比较容差为 $10^{-8}$ pixel。容差内可裁剪到图像边界。}
{\textbf{参考程序：}
\codefile{programs/C24.py}
\textbf{样例输入：}\codefile{snippets/C02.py}
拟合结果为 $x_L=2$、$x_R=7$ pixel，车道中心为 $4.5$ pixel。图像中心为 $(11-1)/2=5$ pixel，偏移为 $+0.5$ pixel。偏移单位为 pixel。没有相机标定和几何映射时，不能把像素偏移称为米。

\textbf{评分：}筛选下半图像并保留原始 $y$ 坐标，3 分；按中心分组并忽略中线，2 分；检查每组具有两个不同 $y$ 值，2 分；拟合并计算底行坐标，3 分；检查坐标范围与左右顺序，2 分；计算中心、偏移并说明单位，3 分。}

\newpage

% C25 S02:15-19 | 综合
\programq{编写 \code{control_metrics(times, speeds, target, band)}，分析一次纵向闭环控制记录，并按多个误差指标评价控制过程。

\textbf{输入：}\code{times} 与 \code{speeds} 是长度相同的非空有限数序列。时间严格递增，单位为 s；速度单位为 m/s。\code{target} 为固定的正有限参考速度，\code{band} 为非负有限误差带。无效输入应引发 \code{ValueError}。

\begin{subquestions}
\item 定义 $e_i=\text{target}-v_i$。计算各样本等权的
\[
\mathrm{RMSE}=\sqrt{\frac{1}{N}\sum_i e_i^2},\qquad E_{\max}=\max_i|e_i|.
\]
（5 分）
\item 计算超调百分比
\[
100\frac{\max(0,\max_i v_i-\text{target})}{\text{target}}.
\]
返回数值不含百分号。（3 分）
\item 定义窗口内稳定起始时刻为最早的 $t_j$：从 $j$ 到最后一个样本均满足 $|e_i|\le\text{band}$，且该后缀至少有 3 个样本。无此时刻则返回 \code{None}。等于误差带时视为满足。返回包含四个指标的字典，并写出样例结果。（7 分）
\end{subquestions}
\codefile{snippets/C03.py}
仅分析给定记录窗口，不推断记录结束后的状态。可用 Python 标准库 \code{math}。}
{\textbf{参考程序：}
\codefile{programs/C25.py}
\textbf{样例输入：}\codefile{snippets/C03.py}
误差为 $[10,-1,0.2,-0.2,-0.1]$ m/s，平方和为 $101.09$。因此 RMSE 为 $\sqrt{101.09/5}\approx4.496$ m/s，最大绝对误差为 $10$ m/s。最大速度为 $11$ m/s，超调百分比为 $10$。下标 1 是最后一个超出误差带的样本，其后 $t=2,3,4$ 的三个样本均在误差带内，故稳定起始时刻为 $2$ s。输出格式见 \code{snippets/C04.py}。

\textbf{评分：}误差符号与 RMSE，3 分；最大绝对误差，2 分；超调计算与非负截断，3 分；检查整个后缀，2 分；至少 3 个样本及最早时刻，3 分；无效输入和无稳定时刻的处理，1 分；样例输出及单位，1 分。}
